% Li Chao Tree

Maintains a set of lines or line segments to query the maximum or minimum value at any integer coordinate $x$ within a fixed range $[L, R]$ in $O(\log(R - L))$.

\begin{lstlisting}
const long long INF = 1e18;

struct Line {
    long long k, m;
    long long operator()(long long x) { return k * x + m; }
};

struct LiChaoTree {
    int n;
    vector<Line> tree;
    LiChaoTree(int n) : n(n), tree(4 * n, {0, -INF}) {} // Set to {0, INF} for Minimum

    void insert(int node, int l, int r, Line newline) {
        int mid = (l + r) >> 1;
        // Change condition to '>' for Minimum
        bool bit_l = newline(l) > tree[node](l);
        bool bit_m = newline(mid) > tree[node](mid);

        if (bit_m) swap(tree[node], newline);
        if (l == r) return;

        if (bit_l != bit_m) insert(node << 1, l, mid, newline);
        else insert(node << 1 | 1, mid + 1, r, newline);
    }

    long long query(int node, int l, int r, int x) {
        if (l == r) return tree[node](x);
        int mid = (l + r) >> 1;
        long long res = tree[node](x);
        if (x <= mid) res = max(res, query(node << 1, l, mid, x)); // Use min() for Minimum
        else res = max(res, query(node << 1 | 1, mid + 1, r, x));   // Use min() for Minimum
        return res;
    }
};
\end{lstlisting}

This default implementation computes the \textbf{MAXIMUM} over the coordinate range $[1, N]$.
\begin{itemize}
    \item \textbf{To compute Minimum:} Change the initialization line dummy to \lstinline|{0, INF}|, replace both line comparisons with \lstinline|>| to \lstinline|<|, and substitute overall node merging functions from \lstinline|max()| to \lstinline|min()|.
\end{itemize}
